\chapter{Progesterone}

\section*{What Is This Test?}
The progesterone test measures the concentration of progesterone in the blood. Progesterone is a C-21 steroid hormone. The corpus luteum of the ovary produces it mainly after ovulation. The placenta produces it during pregnancy. Its principal functions are to prepare and maintain the uterine lining (endometrium) for implantation. It supports early pregnancy. It regulates the menstrual cycle. In the nonpregnant cycling woman, progesterone rises after ovulation. It peaks in the mid-luteal phase. It falls with menses if pregnancy does not occur. A rise above a defined threshold in the luteal phase is the laboratory proof of ovulation. During pregnancy, progesterone maintains the endometrium. It quiets uterine muscle. Serum levels rise progressively with gestational age. The test is used to confirm ovulation. It evaluates luteal phase adequacy in infertility. It assesses abnormal uterine bleeding and amenorrhea. It monitors pregnancy viability, particularly in threatened abortion and suspected ectopic pregnancy. It monitors hormone replacement therapy. Rarely, it evaluates congenital adrenal hyperplasia and progesterone-secreting tumors. In men and postmenopausal women, the adrenal glands produce small amounts of progesterone. It is normally very low in these groups \cite{hoffman2020williams}.

\section*{Alternative Names}
\begin{itemize}
  \item Serum Progesterone
  \item Progesterone Level
  \item P4
  \item Progestin Level
  \item Mid-Luteal Progesterone
  \item Day 21 Progesterone
\end{itemize}
\section*{Principle}
Progesterone is measured by competitive immunoassay methods. The most common are automated chemiluminescent immunoassays (CLIA) or enzyme immunoassays (EIA). Platforms include Roche Cobas, Abbott Architect, Siemens IMMULITE, and Beckman Coulter Access. In a competitive immunoassay, patient progesterone and a labeled progesterone tracer compete. They compete for binding to a fixed amount of anti-progesterone antibody. The signal is inversely proportional to the patient's progesterone concentration. Results are reported in ng/mL or nmol/L (1 ng/mL = 3.18 nmol/L). Liquid chromatography--tandem mass spectrometry (LC-MS/MS) is the reference method. It is preferred in research and in complex clinical settings. It avoids the cross-reactivity and matrix effects of immunoassays \cite{tietz2012textbook}. Progesterone is secreted in a pulsatile pattern. It is cycle-dependent. So the clinical value of a single result depends critically on the timing of collection relative to the menstrual cycle. A mid-luteal progesterone above approximately 3 ng/mL (10 nmol/L) confirms ovulation. Higher values (typically 10--25 ng/mL) reflect adequate corpus luteum function. During pregnancy, levels rise steadily. They are frequently tracked serially to assess viability.

\section*{History}
Progesterone was the second sex steroid hormone identified. It was discovered in 1933--1934 by four independent research groups. George Corner and Willard Allen worked at the University of Rochester. Walter Hohlweg and Karl Slotta worked in Germany. They isolated a substance from the corpus luteum. It maintained pregnancy in ovariectomized animals. The compound was named ``progesterone'' in 1935. The name comes from ``progestational steroidal ketone.'' It was chosen at a meeting of the physiological societies. Adolf Butenandt synthesized it. He received the 1939 Nobel Prize in Chemistry (shared with Leopold Ruzicka) for work on sex hormones. This included progesterone. Synthetic progestins developed from progesterone became the basis of hormonal contraception. Examples are norethindrone and medroxyprogesterone acetate in the 1950s. The modern era of progesterone measurement began with radioimmunoassay (RIA). It used Berson and Yalow's techniques in the 1950s--1960s. Automated immunoassays followed in the 1980s. LC-MS/MS followed in the 2000s. The clinical application of mid-luteal progesterone measurement to confirm ovulation became routine in infertility care in the 1970s--1980s.

\section*{Why This Test Is Ordered}
\begin{itemize}
  \item \textbf{Confirmation of ovulation in infertility evaluation:} A mid-luteal (approximately day 21 of a 28-day cycle, or 7 days before expected menses) progesterone above 3 ng/mL confirms ovulation \cite{speroff2011clinical}
  \item \textbf{Evaluation of luteal phase insufficiency:} In recurrent pregnancy loss or infertility, to determine whether corpus luteum progesterone production is adequate
  \item \textbf{Evaluation of amenorrhea and oligomenorrhea:} To determine whether cycles are ovulatory and to assess the cause of absent or infrequent menses
  \item \textbf{Investigation of abnormal uterine bleeding:} To distinguish ovulatory from anovulatory bleeding
  \item \textbf{Assessment of pregnancy viability:} Serial progesterone measurements in threatened abortion or suspected ectopic pregnancy (a level $<$5 ng/mL is strongly associated with a nonviable pregnancy; levels $>$25 ng/mL are associated with a viable intrauterine pregnancy) \cite{pisarska1999incidence}
  \item \textbf{Monitoring of assisted reproductive technology and progesterone supplementation} in the luteal phase after in vitro fertilization \cite{daya2009luteal}
  \item \textbf{Evaluation of menopausal or perimenopausal status} (low levels) and monitoring of hormone replacement therapy
  \item \textbf{Evaluation of congenital adrenal hyperplasia and adrenal tumors:} Progesterone and 17-hydroxyprogesterone are elevated in some forms; rare progesterone-secreting ovarian or adrenal tumors cause very high levels
\end{itemize}

\section*{Primary vs Secondary Test}
Progesterone is usually a secondary test. It is interpreted within a hormonal panel. This panel includes LH, FSH, estradiol, and, in some settings, testosterone and hCG. It is interpreted with cycle phase. It is the primary laboratory test for confirming ovulation. It is also the primary test for assessing corpus luteum function. In suspected ectopic pregnancy it is a valuable adjunct to serial quantitative hCG. It is not diagnostic by itself.

\section*{How This Test Is Performed}
A venous blood sample is collected in a red-top (serum separator, SST) tube. It is typically 3--5 mL from the antecubital fossa \cite{fischbach2023manual}. The timing of collection is critical. For confirming ovulation, the sample should be drawn approximately 7 days before the expected next menses. This is day 21 of a 28-day cycle. It corresponds to the mid-luteal phase. If cycles are irregular, the draw is timed by ovulation predictor kits, ultrasound monitoring, or repeat sampling. No fasting is required. The sample is analyzed by automated immunoassay. Results are typically available within hours to 1--2 days. The luteinizing hormone surge precedes ovulation by approximately 36 hours. So alternative monitoring (basal body temperature, urinary LH) may be used to time the progesterone draw.

\section*{What Preparation Is Needed}
No fasting is required. The patient should inform the healthcare provider of the date of her last menstrual period. She should also give typical cycle length. She should report any current hormonal medications. These include oral contraceptives, progesterone or progestin therapy, HRT, clomiphene, letrozole, and gonadotropins. She should report whether she may be pregnant. The draw should be scheduled according to the cycle, or timed by the treating physician. This timing must be recorded for accurate interpretation. Morning collection is not specifically required.

\section*{Contraindications and Precautions}
\begin{itemize}
  \item There are no absolute contraindications to the blood draw. Important interpretive cautions:
  \item (1) a single mid-luteal progesterone must be interpreted in relation to cycle day --- an early or late draw may be falsely low;
  \item (2) exogenous progestins and synthetic progestins in contraceptives may cross-react with some immunoassays, while medroxyprogesterone acetate and levonorgestrel generally do not cross-react significantly;
  \item (3) in pregnancy, progesterone does not rise exponentially like hCG, and a single low-normal value is not diagnostic by itself;
  \item (4) in ectopic pregnancy, progesterone is frequently low ($<$10 ng/mL), but up to 10--15\% of viable intrauterine pregnancies also have low values;
  \item (5) 17-hydroxyprogesterone is structurally similar and may cross-react in some assays, relevant in congenital adrenal hyperplasia;
  \item (6) because progesterone is elevated by ovulation, results cannot be interpreted in anovulatory or postmenopausal patients without reference to the clinical context. Certain rare ovarian and adrenal tumors secrete progesterone and cause markedly elevated levels.
\end{itemize}
\section*{Risks}
Risks are those of routine venipuncture. They include mild pain, bruising, and small hematoma at the collection site. They also include lightheadedness or vasovagal syncope. Extremely rare risks are infection or nerve injury.

\section*{What Happens After the Test}
Results are typically available within 1--2 days. In infertility evaluation, a mid-luteal progesterone above 3 ng/mL confirms ovulation. Values below this suggest anovulation or mis-timed sampling. The test is often repeated in a subsequent cycle or timed by ultrasound. In pregnancy, serial progesterone and serial quantitative hCG measurements are used. They discriminate viable intrauterine pregnancy from ectopic pregnancy or threatened abortion. A level below 5 ng/mL is strongly associated with a nonviable pregnancy. Levels above 25 ng/mL are reassuring. Progesterone is required to maintain early pregnancy. So patients with low levels and a desired pregnancy may be treated with progesterone supplementation. This can be vaginal, oral, or intramuscular. This intervention has been shown to reduce miscarriage in women with early pregnancy bleeding and a previous miscarriage \cite{coomarasamy2019randomized}. Patients on luteal-phase support after IVF have progesterone monitored. Dosing is adjusted accordingly. The healthcare provider will always interpret the progesterone result together with the cycle day, hCG, ultrasound findings, and other reproductive hormones.

\section*{Understanding Results}

\subsection*{Below Normal (in the Luteal Phase or Pregnancy)}
\begin{itemize}
  \item \textbf{Anovulation:} Low progesterone throughout the cycle, with no mid-luteal rise --- seen in anovulatory bleeding, polycystic ovary syndrome (PCOS), perimenopause, and hypothalamic or hyperprolactinemic amenorrhea
  \item \textbf{Luteal phase insufficiency:} Mid-luteal progesterone below approximately 3--10 ng/mL, associated with infertility and recurrent pregnancy loss
  \item \textbf{Nonviable pregnancy or ectopic pregnancy:} Low or falling progesterone in pregnancy (commonly $<$5--10 ng/mL) with abnormal serial hCG
  \item \textbf{Postmenopausal state:} Low progesterone (<1 ng/mL) is expected; elevated levels suggest exogenous therapy or a rare progesterone-secreting tumor
  \item \textbf{Ovarian failure and menopause:} Absent luteal phase and low progesterone
\end{itemize}

\subsection*{Above Normal}
\begin{itemize}
  \item \textbf{Mid-luteal phase (ovulation):} Normal physiological peak of 3--25 ng/mL
  \item \textbf{Pregnancy:} Progressive rise with gestation, reaching 100 ng/mL and higher in the third trimester
  \item \textbf{Molar pregnancy:} Elevated progesterone and markedly elevated hCG
  \item \textbf{Exogenous progesterone or progestin therapy:} Vaginal, oral, or intramuscular supplementation raises measured levels
  \item \textbf{Congenital adrenal hyperplasia:} Elevated precursors including 17-hydroxyprogesterone may cross-react
  \item \textbf{Rare ovarian or adrenal progesterone-secreting tumors:} Very high levels with virilization or Cushing syndrome features
\end{itemize}

\section*{Normal Results}
\begin{itemize}
  \item \textbf{Follicular phase:} 0.1--1.0 ng/mL (0.3--3.2 nmol/L)
  \item \textbf{Mid-luteal phase (ovulatory):} 3.0--25 ng/mL (10--80 nmol/L); a value $>$3 ng/mL confirms ovulation, and $>$10 ng/mL reflects adequate corpus luteum function
  \item \textbf{Postmenopausal women:} $<$0.5--1.0 ng/mL
  \item \textbf{Men:} $<$1.0 ng/mL (reference ranges vary by laboratory)
  \item \textbf{Pregnancy (approximate, rising with gestation):}
  \begin{itemize}
    \item First trimester: 10--30 ng/mL (rising)
    \item Second trimester: 20--90 ng/mL
    \item Third trimester: 100--300 ng/mL
  \end{itemize}
\end{itemize}

\section*{Follow-up Tests}
\begin{itemize}
  \item \textbf{Quantitative hCG:} For pregnancy evaluation, especially when ectopic pregnancy or threatened abortion is suspected
  \item \textbf{Luteinizing hormone and follicle-stimulating hormone (LH/FSH):} To evaluate the menstrual cycle and ovulatory function
  \item \textbf{Estradiol:} With LH/FSH to characterize the menstrual cycle and ovarian reserve
  \item \textbf{Testosterone and dehydroepiandrosterone sulfate (DHEA-S):} In hyperandrogenic states such as PCOS and adrenal disorders
  \item \textbf{17-Hydroxyprogesterone:} To diagnose congenital adrenal hyperplasia when precursors are suspected
  \item \textbf{Thyroid-stimulating hormone (TSH) and prolactin:} Common causes of ovulatory dysfunction
  \item \textbf{Pelvic ultrasound:} To assess the endometrium, corpus luteum, follicle size, and adnexal masses
  \item \textbf{Endometrial biopsy or luteal-phase timed biopsy:} In selected luteal phase assessment protocols
  \item \textbf{Hysteroscopy and hysterosalpingography:} In the structural evaluation of infertility
\end{itemize}

\section*{References}
\bibliographystyle{unsrtnat}
\bibliography{references}

\clearpage
\section*{Regulatory Status and Authorization Date}
This chapter describes a test category, not a specific commercial product. FDA, CDSCO, EU IVDR, and MHRA approval or registration dates therefore cannot be assigned without the assay manufacturer, product name, instrument, and intended use. For laboratory-developed tests, an individual agency approval date may not exist. See the Preface for the interpretation of regulatory status.

\clearpage
